% Part: first-order-logic
% Chapter: sequent-calculus

\documentclass[../../../include/open-logic-chapter]{subfiles}

\begin{document}

\iftag{FOL}
      {\olchapter{fol}{seq}{Kalkulus Sekuen}}
      {\olchapter{pl}{seq}{Kalkulus Sekuen}}

\begin{editorial}
  Bab iki nyajekake kalkulus sekuen baku LK anggitane Gentzen kanggo
  logika klasik tataran kapisan. Bab iki isih mbutuhake luwih akeh
  tuladha lan gladhen. Kanggo ngemot utawa nyingkirake materi kang
  gegayutan karo kalkulus sekuen minangka sistem bukti, gunakna tag
  “prfLK”.
\end{editorial}

\olimport{rules-and-proofs}

\olimport{propositional-rules}

\iftag{FOL}{%
  \olimport{quantifier-rules}
}{}

\olimport{structural-rules}

\olimport{derivations}

\olimport{proving-things}

\iftag{FOL}{%
  \olimport{proving-things-quant}
}{}

\olimport{proof-theoretic-notions}

\olimport{provability-consistency}

\olimport{provability-propositional}

\iftag{FOL}{%
  \olimport{provability-quantifiers}
}{}

\olimport{soundness}

\iftag{FOL}{%
  \olimport{identity}
  \olimport{soundness-identity}
}{}

\OLEndChapterHook

\end{document}
